\section{Electronics and power}
\label{sec:electronics}

\subsection{Architecture}
\label{sec:architecture}

Figure~\ref{fig:electronics} shows the wiring. A Particle Photon~2~\cite{photon2} runs the robot. It replaced the kit's ESP32 board, which failed during bench wiring for reasons we could not determine. The kit's motor shield drives its motors through a co-processor commanded over a serial link; we replaced it with a SparkFun TB6612FNG dual motor driver~\cite{sparkfun} we had on hand, controlled directly from the Photon's pins and with its motor supply wired straight to the battery. An MPU6050~\cite{mpu6050} gyro and accelerometer on the I$^2$C bus does all the sensing, read by firmware written mostly with Claude Code~\cite{claude}. Power comes from the supplied 2S pack of 18650 cells \tbd{confirm capacity: 2,500 or 3,000 mAh}. A judge starts the run with a single switch, which the Photon reads directly.

\begin{figure}[htbp]
  \centering
  \definecolor{packred}{HTML}{E8402F}\definecolor{fiveorange}{HTML}{E08A00}\definecolor{logicblue}{HTML}{1F6FB4}
  \begin{tikzpicture}[font=\small,
      box/.style={draw=black,line width=0.8pt,rounded corners=2pt,fill=black!4,align=center,minimum height=1.3cm,text width=2.5cm},
      pack/.style={line width=1.6pt,draw=packred,-{Latex[length=2.2mm]}},
      five/.style={line width=1.6pt,draw=fiveorange,-{Latex[length=2.2mm]}},
      logic/.style={line width=1.2pt,draw=logicblue,-{Latex[length=2.2mm]}},
      sig/.style={line width=1.0pt,draw=black,-{Latex[length=2.2mm]}}]
    \node[box,fill=black!2,text width=2.3cm] (pack) at (0,0) {\textbf{2S pack}\\two 18650 cells};
    \node[box] (reg) at (3.4,1.6) {\textbf{7805}\\regulator};
    \node[box] (sw) at (3.4,-0.6) {\textbf{Start switch}};
    \node[box,minimum height=4.2cm,text width=2.4cm,fill=black!2] (mcu) at (8.4,0.2) {\textbf{Photon 2}\\controller\\(\qty{3.3}{V} logic)};
    \node[box] (imu) at (13.2,1.6) {\textbf{MPU6050}\\gyro and\\accelerometer};
    \node[box,text width=2.9cm] (drv) at (13.2,-1.6) {\textbf{TB6612FNG}\\dual motor driver};
    \node[box,text width=1.4cm,minimum height=1cm] (ml) at (12.3,-3.8) {\textbf{Left}\\motor};
    \node[box,text width=1.4cm,minimum height=1cm] (mr) at (14.1,-3.8) {\textbf{Right}\\motor};
    \draw[pack] (pack.north) |- node[pos=0.75,above,text=packred] {pack} (reg.west);
    \draw[five] (reg.east) -- node[above,text=fiveorange] {\qty{5}{V} (VUSB)} (reg.east -| mcu.west);
    \draw[sig] (sw.east) -- node[above] {digital in} (sw.east -| mcu.west);
    \draw[logic,{Latex[length=2.2mm]}-{Latex[length=2.2mm]}] (imu.west -| mcu.east) -- node[above,text=logicblue] {I$^2$C} (imu.west);
    \draw[logic] ([yshift=0.3cm]drv.west -| mcu.east) -- node[above,text=logicblue] {control} ([yshift=0.3cm]drv.west);
    \draw[pack] (pack.south) |- node[pos=0.75,below,text=packred] {VM: pack voltage (motor supply)} (11.0,-2.7) |- ([yshift=-0.35cm]drv.west);
    \draw[sig] (ml.north |- drv.south) -- (ml.north);
    \draw[sig] (mr.north |- drv.south) -- (mr.north);
  \end{tikzpicture}
  \caption{Electronics block diagram. The 2S pack feeds the motor driver directly and the Photon~2 through the 7805 regulator. The Photon reads the MPU6050 over I$^2$C and drives the TB6612FNG.}
  \label{fig:electronics}
\end{figure}

The TB6612FNG switches the motors with MOSFETs and loses about \qty{0.5}{V} at \qty{1}{A}, against about \qty{2}{V} for the bipolar L298 used on many hobby driver boards, so more of the pack voltage reaches the motors~\cite{tb6612,l298}.

The Photon takes \qty{5}{V} on its VUSB pin and makes its own \qty{3.3}{V} for logic, so every sensor and signal line runs at \qty{3.3}{V}. A 7805 linear regulator makes the \qty{5}{V} from the pack. It needs about \qty{2}{V} of headroom, so it holds \qty{5}{V} only while the pack stays above about \qty{7}{V}; a 2S pack runs from \qty{8.4}{V} full to about \qty{6.4}{V} empty~\cite{cell}. We bench-tested the regulator with the capacitors we had, two \qty{100}{nF} in parallel on the input and one \qty{100}{nF} on the output, below the \qty{0.33}{\micro F} input capacitor the datasheet calls for~\cite{reg7805}, and measured \qty{5.05}{V} out.

\subsection{Resets from a loose cell}
\label{sec:resets}

During bench testing the Photon sometimes restarted when the robot was jolted. Stalling both motors never caused a restart, but snapping the robot back and forth by hand with the motors idle did, which ruled out a supply dip from motor current. The cells sit in a plastic holder with spring contacts, and the jolt let a cell compress its spring far enough to lose contact for an instant. Fastening the holder so it could not snap against its stop ended the restarts \tbd{confirm}.
